%==============================
\section{Case study}
\label{sec:ch05_case_study}

A representative scenario for this chapter combines the three technical themes
of communication-aware planning, intermittent connectivity, and online temporal
requests within a single exploration-and-service mission. Consider a robot
fleet deployed in a partially known workspace, where the team must continue
exploration, satisfy latency bounds on information propagation, and remain able
to accommodate operator-issued requests during execution. At the outset, the
robots share only a partial map and therefore cannot rely on a fixed
communication infrastructure or on complete prior knowledge of the relevant task
locations. The communication topology must instead be planned together with the
team behavior.

The mission begins with the communication-and-collaboration mechanism of
Section~\ref{ch05:sec:ch05_interm_team}. The robots execute their current local
tasks while scheduling communication events that are informative enough to keep
the distributed plans mutually consistent. As exploration proceeds, the setting
transitions naturally to the intermittent-communication regime of
Section~\ref{ch05:sec:intermittent_comm}. Rather than enforcing persistent
connectivity, the team maintains bounded-delay information fusion through the
embedded joint-wheel topology and the detach-and-rejoin strategy. This allows
subgroups to move away from the backbone temporarily to inspect frontiers or
service nearby objectives, while still guaranteeing that accumulated data and
plan updates are reintegrated after scheduled meetings.

Now suppose that an operator introduces a temporal request during execution,
for example by requiring a particular subgroup to inspect a newly designated
region, deliver a payload, or provide local assistance under a precedence
constraint. The request-fulfillment mechanism of
Section~\ref{ch05:sec:ch05_human} evaluates whether the current wheel backbone
already supports the new requirement or whether a temporary request-specific
subtopology must be instantiated. If the request is topology-preserving, it is
absorbed into the existing structure through an embedding update. If it is
topology-changing, a compatible subgroup is detached, executes the request, and
later rejoins the backbone so that exploration and communication maintenance can
continue. This case-study synthesis clarifies the main message of the chapter:
under partial knowledge and sparse communication, coordination is governed not
only by what the robots do, but also by how and when the information needed for
those actions is exchanged.

%==============================
\section{Summary}
\label{sec:ch05_summary}


This chapter examined multi-robot coordination under partial knowledge
and limited communication, with emphasis on how communication
constraints should be treated as part of the planning problem rather
than as a passive execution condition. The chapter first considered the
joint design of communication and collaboration, where robot actions
and information exchange must be coordinated simultaneously in order to
maintain task feasibility. It then developed an exploration framework
under intermittent communication, in which the operator and the robot
fleet preserve bounded-delay information fusion through an embedded
joint-wheel communication topology while continuing frontier-based
exploration in an initially unknown workspace. On this basis, the
chapter introduced a topology-adaptive coordination layer for online
temporal requests, where the wheel backbone serves as the nominal
communication structure and is modified only when a request requires a
temporary reconfiguration of the embedded topology.

Taken together, the developments in
Sections~\ref{ch05:sec:ch05_interm_team}--
\ref{ch05:sec:ch05_human} provide a unified view of multi-robot
coordination under sparse and delayed communication. Communication
maintenance, exploration, and request fulfillment are integrated within
a common planning and execution architecture, rather than treated as
separate modules. The resulting framework preserves latency-bounded
information flow whenever possible, supports online adaptation to newly
discovered workspace information and operator requests, and allows
temporary request-specific subtopologies to be instantiated and later
merged back into the wheel backbone. In this way,
Chapter~\ref{chap:communication} extends the earlier coordination
frameworks of Chapters~\ref{chap:bottom-up} and~\ref{chap:top-down} to
settings in which communication is intermittent, information is
incomplete, and topology adaptation is required for continued task
fulfillment.
